\documentclass{article}
\usepackage[T1]{fontenc}
\usepackage{lmodern}
\usepackage{graphicx}
\usepackage{amsmath,amssymb}
\usepackage[a4paper,margin=2cm]{geometry}
\usepackage[absolute,overlay]{textpos}
\setlength{\TPHorizModule}{1mm}
\setlength{\TPVertModule}{1mm}
\usepackage[indonesian]{babel}
\usepackage{hyperref}
\setlength{\emergencystretch}{2em}
% unit: Halaman 6

\begin{document}

% === Halaman 6 (python/native) ===

• Sebagai contoh, misalkan p = 7, maka a = 3 adalah akar primitif dari 7 karena

31 (mod 7) = 3;

32 (mod 7) = 2;

33 (mod 7) = 6

34 (mod 7) = 4;

35 (mod 7) = 5;

36 (mod 7) = 1

• Jadi, semua perpangkatan dari 3 menghasilkan nilai-nilai yang berbeda (3, 2, 6, 4, 

5, 1), semua bilangan di dalam modulus 7 terjadi satu kali. 

• Perpangkatan berikutnya, 37 (mod 7), 38 (mod 7), …, akan kembali berulang 

menghasilkan nilai-nilai tersebut.  Panjang satu siklus tidak lebih dari 7 – 1 = 6.

• Secara umum, untuk p bilangan prima, maka panjang siklus tidak lebih dari p – 1.  

6

\end{document}
